\documentclass[12pt,a4paper]{article}

% ============================================================================
% Packages
% ============================================================================
\usepackage[utf-8]{inputenc}
\usepackage[T1]{fontenc}
\usepackage{geometry}
\usepackage{amsmath}
\usepackage{amssymb}
\usepackage{graphicx}
\usepackage{hyperref}
\usepackage{listings}
\usepackage{xcolor}
\usepackage{fancyhdr}
\usepackage{booktabs}
\usepackage{caption}
\usepackage{subcaption}
\usepackage{natbib}
\usepackage{doi}

% ============================================================================
% Geometry and margins
% ============================================================================
\geometry{
  left=2.5cm,
  right=2.5cm,
  top=2.5cm,
  bottom=2.5cm,
  headsep=0.5cm,
  footskip=1cm
}

% ============================================================================
% Code listings configuration
% ============================================================================
\lstset{
  language=Python,
  basicstyle=\ttfamily\small,
  commentstyle=\color{gray},
  keywordstyle=\color{blue}\bfseries,
  stringstyle=\color{red},
  breaklines=true,
  breakatwhitespace=true,
  prebreak=\raisebox{0ex}[0ex][0ex]{\ensuremath{\hookleftarrow}},
  postbreak=\raisebox{0ex}[0ex][0ex]{\ensuremath{\hookrightarrow\space}},
  showstringspaces=false,
  lineskip=2pt,
  frame=single,
  frameround=tttt,
  framexleftmargin=3pt,
  framexrightmargin=3pt,
  framextopmargin=4pt,
  framexbottommargin=4pt,
  backgroundcolor=\color{lightgray!10},
  numbers=left,
  numberstyle=\tiny\color{gray},
  stepnumber=1,
}

% ============================================================================
% Header and footer
% ============================================================================
\pagestyle{fancy}
\fancyhf{}
\rhead{DDM501 Lab 1}
\lhead{Movie Rating Prediction ML Product}
\cfoot{\thepage}

% ============================================================================
% Title configuration
% ============================================================================
\title{%
  \Large\textbf{Movie Rating Prediction:}\\
  \large A Machine Learning Product from Training to Deployment\\
  \large DDM501 Lab 1: First ML Product
}
\author{%
  Master of Science in Engineering (MSE)\\
  \vspace{0.3cm}
  Data Science and Engineering Track
}
\date{April 2026}

% ============================================================================
% Hyperref configuration
% ============================================================================
\hypersetup{
  colorlinks=true,
  linkcolor=blue,
  filecolor=blue,
  urlcolor=blue,
  citecolor=blue,
  pdftitle={Movie Rating Prediction ML Product},
  pdfauthor={DDM501 Lab Course},
}

% ============================================================================
% Custom commands
% ============================================================================
\newcommand{\code}[1]{\texttt{#1}}
\newcommand{\citationneeded}{\textsuperscript{[citation needed]}}

% ============================================================================
% Document
% ============================================================================
\begin{document}

% Title page
\maketitle

% ============================================================================
% Abstract
% ============================================================================
\begin{abstract}
  This report documents the design, implementation, and deployment of a movie rating prediction 
  system as a complete machine learning product. The system leverages collaborative filtering 
  via Singular Value Decomposition (SVD) to predict user ratings for movies, trained on the 
  MovieLens 100K dataset. The machine learning model is served through a RESTful API built 
  with FastAPI, containerized using Docker, and validated through comprehensive pytest-based 
  testing. This work demonstrates the full lifecycle of a production ML system: from data 
  preparation and model training to API design, deployment, and operational validation. The system 
  achieves robust performance with graceful handling of edge cases, and serves as a practical 
  reference implementation for ML product development.
\end{abstract}

% ============================================================================
% Table of contents
% ============================================================================
\newpage
\tableofcontents

\newpage

% ============================================================================
% Sections
% ============================================================================
\section{Introduction}

Machine learning has fundamentally transformed software development, but deploying ML systems to production 
remains a significant engineering challenge. The gap between successful model development and reliable 
production deployment is often underestimated in academic settings. This report documents a complete, 
end-to-end machine learning product: a movie rating prediction system that moves beyond model training 
to address the full lifecycle of an ML system.

\subsection{Motivation}

The movie recommendation domain serves as an excellent case study for several reasons:

\begin{enumerate}
  \item \textbf{Clear business value}: Movie rating predictions directly support recommendation engines, 
  which are critical for user engagement and platform success.
  
  \item \textbf{Rich collaborative patterns}: User-movie interactions exhibit strong collaborative signal, 
  where similar users rate similar movies similarly. This structure makes collaborative filtering effective.
  
  \item \textbf{Tractable scale}: The MovieLens 100K dataset is large enough to demonstrate real-world 
  challenges (sparsity, computation) while remaining manageable for classroom implementation.
  
  \item \textbf{Accessible baselines}: Collaborative filtering is well-understood and has established 
  theoretical foundations, allowing focus on engineering best practices rather than algorithmic novelty.
\end{enumerate}

\subsection{Scope of Work}

This project addresses three major engineering domains:

\textbf{Machine Learning:} We train a Singular Value Decomposition (SVD) model on the MovieLens 100K dataset 
to perform collaborative filtering. The model learns latent factor representations of users and movies, 
enabling rating predictions for arbitrary user-movie pairs.

\textbf{API and Systems Design:} We implement a RESTful API using FastAPI, providing clean interfaces 
for health monitoring and prediction serving. The API enforces input validation, handles errors gracefully, 
and maintains separation of concerns between model logic and request handling.

\textbf{Operations and Deployment:} We containerize the system using Docker and Docker Compose, ensuring 
portable and reproducible deployment. Health checks, environment configuration, and volume management 
demonstrate production-ready practices.

\subsection{Key Contributions}

This work makes the following contributions:

\begin{enumerate}
  \item A complete, working implementation of an ML system suitable for production-like constraints.
  
  \item Practical demonstration of API design patterns (validation, error handling, versioning) for ML systems.
  
  \item Comprehensive testing strategy that balances functional validation with realistic edge case handling.
  
  \item Dockerized deployment with health checks, supporting both local and containerized execution.
  
  \item Documentation and code structure that prioritizes maintainability and team collaboration.
\end{enumerate}

\subsection{Report Structure}

The remainder of this report is organized as follows:

\begin{itemize}
  \item \textbf{Section~\ref{sec:problem}}: Formal problem definition and rating prediction task formulation.
  
  \item \textbf{Section~\ref{sec:dataset}}: MovieLens 100K dataset characteristics, statistics, and preprocessing.
  
  \item \textbf{Section~\ref{sec:methodology}}: Theoretical foundations of collaborative filtering and matrix factorization.
  
  \item \textbf{Section~\ref{sec:model}}: Mathematical treatment of SVD and the Surprise library implementation.
  
  \item \textbf{Section~\ref{sec:architecture}}: System architecture covering training, artifact management, and serving.
  
  \item \textbf{Section~\ref{sec:api}}: RESTful API design, schema definitions, and error handling.
  
  \item \textbf{Section~\ref{sec:experiments}}: Training results, model evaluation, and prediction behavior analysis.
  
  \item \textbf{Section~\ref{sec:deployment}}: Docker containerization and deployment configuration.
  
  \item \textbf{Section~\ref{sec:testing}}: Testing strategy, test cases, and validation approach.
  
  \item \textbf{Section~\ref{sec:limitations}}: System limitations, failure modes, and future improvements.
  
  \item \textbf{Section~\ref{sec:conclusion}}: Summary and directions for future work.
\end{itemize}

\vspace{0.5cm}

This report emphasizes the engineering and operational aspects of ML systems, recognizing that 
successful ML in production requires far more than accurate model training.

\section{Problem Definition}
\label{sec:problem}

\subsection{Rating Prediction Task}

The core problem addressed by this system is \textit{rating prediction}: given a user $u$ and an item 
(movie) $i$, predict the numerical rating $r_{u,i}$ that the user would assign to the item. Formally:

\begin{equation}
  f: (u, i) \mapsto \hat{r}_{u,i} \in [1, 5]
\end{equation}

where:
\begin{itemize}
  \item $u \in \mathcal{U}$ is a user identifier from the set of all users.
  \item $i \in \mathcal{I}$ is an item (movie) identifier from the set of all items.
  \item $\hat{r}_{u,i}$ is the predicted rating, constrained to the interval $[1.0, 5.0]$.
\end{itemize}

The prediction is made using historical rating data $\mathcal{R} = \{(u_j, i_j, r_{u_j, i_j})_{j=1}^{|\mathcal{R}|}\}$, 
representing past observations of users' ratings on movies.

\subsection{Setting and Assumptions}

We operate under the following setting:

\textbf{Implicit Feedback Context:} While the MovieLens 100K dataset contains explicit numerical ratings, 
the system must handle both known and unknown user-item pairs. For unknown pairs (users or movies not in the 
training set), the model is expected to produce reasonable estimates based on learned latent factors.

\textbf{Cold Start Simplifications:} We do not implement sophisticated cold-start strategies. Instead, we 
rely on the base model's ability to extrapolate from latent factor space. This is a practical limitation 
addressed in Section~\ref{sec:limitations}.

\textbf{Static Model Assumption:} The model is trained once and served in production without online learning. 
This is typical for batch-trained systems and simplifies operational deployed.

\textbf{Single Prediction Paradigm:} Beyond single-prediction requests, the API also supports batch 
predictions---processing multiple $(u, i)$ pairs in a single request. Both modes share the same underlying 
prediction logic.

\subsection{Formal Data Representation}

The rating data is typically represented as a sparse matrix $R \in \mathbb{R}^{|\mathcal{U}| \times |\mathcal{I}|}$, 
where entry $R_{u,i}$ is the rating given by user $u$ to movie $i$, or undefined if the rating was never 
provided. This sparsity is fundamental to the problem: real-world user-item interactions are sparse, 
and any effective recommender system must handle this sparsity gracefully.

\subsection{Success Criteria}

The system is evaluated against the following criteria:

\begin{enumerate}
  \item \textbf{Correctness}: Predictions are real-valued and fall within $[1.0, 5.0]$.
  
  \item \textbf{Reliability}: The system handles edge cases (unknown users, missing data) without crashing.
  
  \item \textbf{Performance}: Prediction latency is acceptable (sub-second) for serving via REST API.
  
  \item \textbf{Usability}: The API is well-documented, validates input, and returns meaningful error messages.
  
  \item \textbf{Operability}: The system can be deployed, monitored, and scaled using standard DevOps practices.
\end{enumerate}

These criteria span model accuracy, system robustness, and operational concerns---reflecting the reality 
that production ML systems are judged on far more than raw prediction accuracy.

\section{Dataset}
\label{sec:dataset}

\subsection{MovieLens 100K Overview}

The MovieLens 100K dataset is a widely-used benchmark in recommendation systems research. It was collected 
by the GroupLens research group at the University of Minnesota over a seven-month period (Sept 19th, 1997 -- 
April 22nd, 1998). The dataset consists of approximately 100,000 ratings provided by around 1,000 users on 
1,700 movies.

\textbf{Data Distribution:}
\begin{itemize}
  \item \textbf{Number of users}: 943
  \item \textbf{Number of items (movies)}: 1,682
  \item \textbf{Number of ratings}: 100,000
  \item \textbf{Rating scale}: 1--5 (integer values)
  \item \textbf{Sparsity}: $\frac{100,000}{943 \times 1,682} \approx 6.3\%$ (about 93.7\% of the matrix is missing)
\end{itemize}

The high sparsity is typical of real-world recommendation problems and is central to the challenge: 
predicting ratings for user-item pairs that have never been observed.

\subsection{Data Format and Preprocessing}

The MovieLens 100K dataset is provided in a tabular format. Each row represents a single rating observation, 
with fields:
\begin{itemize}
  \item User ID: integer identifier (1--943)
  \item Item ID: integer identifier (1--1,682)
  \item Rating: integer in $\{1, 2, 3, 4, 5\}$
  \item Timestamp: Unix timestamp of the rating event
\end{itemize}

The \code{Surprise} library handles dataset loading and preprocessing automatically. Specifically:

\begin{enumerate}
  \item \textbf{String conversion}: User and item IDs are converted to strings for API compatibility.
  \item \textbf{Train-test splits}: The dataset provides standard cross-validation folds (u1--u5) and 
  alternative train-test splits (ua, ub) for benchmarking.
  \item \textbf{No additional preprocessing}: Ratings are used as-is. No normalization, outlier removal, 
  or feature engineering is applied in the base implementation.
\end{enumerate}

\subsection{Statistical Characteristics}

Analysis of the MovieLens 100K dataset reveals:

\textbf{Rating Distribution:}
The ratings follow a roughly normal distribution with slight skew toward positive ratings. The mean rating 
is approximately 3.5, indicating users tend to rate items they watch (selection bias).

\textbf{User Activity:}
Users vary dramatically in activity level. Some users have provided hundreds of ratings, while others 
contributed only a handful. This heterogeneity impacts model training: high-activity users provide more 
signal, while low-activity (cold-start) users are harder to model.

\textbf{Item Popularity:}
Similarly, movies vary in the number of ratings they receive. Popular movies have many ratings, enabling 
reliable factor estimation, while niche movies have few ratings and are prone to cold-start effects.

\textbf{Temporal Characteristics:}
Though the dataset is small by modern standards and temporally localized to 1997--1998, the timestamps allow 
for temporal analysis. For this project, temporal information is not utilized, though it could support 
cold-start improvements or drift detection in production.

\subsection{Data Assumptions and Limitations}

Important assumptions about the data:

\begin{enumerate}
  \item \textbf{Ratings represent user preferences}: We assume that user ratings reflect their true preferences 
  (no gaming or manipulation).
  
  \item \textbf{Collaborative signal exists}: We assume that similar users rate similar movies similarly. 
  This underpins the collaborative filtering approach.
  
  \item \textbf{Data is clean}: The dataset is curated and does not require extensive cleaning. In production, 
  data validation and quality checks would be essential.
  
  \item \textbf{Static snapshot}: The dataset is treated as a static snapshot. In production, the training 
  data would be continuously refreshed.
  
  \item \textbf{No side information}: We use only ratings. User demographics, movie metadata, and tags are 
  not utilized (though they could improve predictions).
\end{enumerate}

\subsection{Sparsity and Its Implications}

The sparsity of the rating matrix is the central challenge in recommendation systems:

\[
  \text{Sparsity} = 1 - \frac{|\mathcal{R}|}{|\mathcal{U}| \times |\mathcal{I}|} = 1 - \frac{100,000}{943 \times 1,682} \approx 93.7\%
\]

Sparsity implies that:
\begin{itemize}
  \item Most user-item pairs have no observed rating.
  \item The rating matrix cannot be inverted or directly analyzed using standard matrix methods.
  \item Extrapolation from limited information (latent factors) is necessary.
  \item Cold-start (new users or items with few ratings) is a persistent problem.
\end{itemize}

Collaborative filtering through matrix factorization addresses sparsity by learning low-rank latent 
representations (Section~\ref{sec:methodology}).

\section{Methodology}
\label{sec:methodology}

\subsection{Collaborative Filtering}

Collaborative filtering is a family of recommendation techniques based on the premise that users who have 
agreed in the past will agree in the future. The core insight is that user preferences are often correlated: 
if user $u_1$ and user $u_2$ have rated several movies similarly, then we can use movies rated by $u_1$ to 
recommend movies to $u_2$.

Formally, the assumption is:

\[
  \text{If } \text{sim}(u_i, u_j) \text{ is high, then unknown ratings by user } u_i 
  \text{ can be inferred from user } u_j \text{'s ratings.}
\]

\subsection{Memory-Based vs.~Model-Based Approaches}

Collaborative filtering is typically categorized into two classes:

\textbf{Memory-Based Approaches:}
\begin{itemize}
  \item Store the full rating matrix in memory.
  \item At prediction time, compute similarity between the target user and all other users.
  \item Predict by averaging ratings from similar users (weighted by similarity).
  \item Advantages: Interpretable, fast for small datasets, requires no training.
  \item Disadvantages: Scales poorly ($O(|\mathcal{U}|^2)$ per prediction), cannot handle sparsity well.
\end{itemize}

\textbf{Model-Based Approaches:}
\begin{itemize}
  \item Learn a compact representation of users and items.
  \item During training, optimize parameters to reconstruct observed ratings.
  \item Predictions are computed using the learned model, not raw data.
  \item Advantages: Scalable ($O(k)$ per prediction, where $k$ is the latent dimension), handles sparsity 
  better, supports extrapolation.
  \item Disadvantages: Requires training, less interpretable, prone to overfitting.
\end{itemize}

This project employs a model-based approach via matrix factorization.

\subsection{Matrix Factorization and SVD}

Matrix factorization is a model-based approach that decomposes the rating matrix $R$ into a product of 
lower-rank matrices:

\[
  R \approx U \Sigma V^T
\]

where:
\begin{itemize}
  \item $U \in \mathbb{R}^{|\mathcal{U}| \times k}$ is the user factor matrix (low-rank approximation of user preferences).
  \item $\Sigma \in \mathbb{R}^{k \times k}$ is a diagonal matrix of singular values (importance of each latent factor).
  \item $V^T \in \mathbb{R}^{k \times |\mathcal{I}|}$ is the transpose of the item factor matrix (low-rank approximation of item characteristics).
  \item $k$ is the rank or number of latent factors (typically $k \ll \min(|\mathcal{U}|, |\mathcal{I}|)$).
\end{itemize}

The prediction for user $u$ on item $i$ is:

\[
  \hat{r}_{u,i} = U_u \Sigma V_i^T
\]

where $U_u$ is the $u$-th row of $U$ and $V_i$ is the $i$-th column of $V$ (or equivalently, the $i$-th row of $V^T$).

\subsection{Rationale for Matrix Factorization}

Matrix factorization is chosen for this project because:

\begin{enumerate}
  \item \textbf{Handles sparsity}: By learning latent factors, the model can predict ratings for 
  unobserved user-item pairs.
  
  \item \textbf{Computational efficiency}: Inference is $O(k)$, enabling real-time predictions.
  
  \item \textbf{Scalability}: Training can be parallelized; the model size grows with $k$, not dataset size.
  
  \item \textbf{Theoretical foundation}: SVD has well-understood mathematical properties and convergence guarantees 
  under standard assumptions.
  
  \item \textbf{Proven empirical performance}: Matrix factorization achieves state-of-the-art results on 
  benchmark datasets when properly tuned.
  
  \item \textbf{Library support}: The \code{Surprise} library provides robust implementations.
\end{enumerate}

\subsection{Optimization Objective}

During training, the model minimizes a loss function that encourages the reconstructed matrix to match 
observed ratings while avoiding overfitting. A standard formulation is:

\[
  \min_{U, V} \sum_{(u,i) \in \mathcal{R}} (r_{u,i} - U_u V_i^T)^2 + \lambda_U \|U\|_F^2 + \lambda_V \|V\|_F^2
\]

where:
\begin{itemize}
  \item The first term is the reconstruction error (mean squared error on observed ratings).
  \item The second and third terms are $L_2$ regularization penalties on user and item factors.
  \item $\|.\|_F$ denotes the Frobenius norm.
  \item $\lambda_U$ and $\lambda_V$ are regularization hyperparameters controlling overfitting.
\end{itemize}

This objective is optimized using stochastic gradient descent (SGD) or alternating least squares (ALS). 
The \code{Surprise} library uses SGD by default, with configurable learning rates and regularization strengths.

\subsection{Why Not Deep Learning?}

While neural networks (autoencoders, NCF) are state-of-the-art for recommendation, this project uses SVD 
because:

\begin{enumerate}
  \item \textbf{Interpretability}: SVD factors have meaningful structure and can be analyzed.
  \item \textbf{Training efficiency}: SVD converges quickly on modest hardware.
  \item \textbf{Simplicity}: SVD is well-understood and implemented reliably.
  \item \textbf{Laboratory setting}: This project prioritizes breadth (full ML product cycle) over 
  cutting-edge model performance.
\end{enumerate}

In production, migrating to neural approaches would be a straightforward extension.

\section{Model: SVD Collaborative Filtering}
\label{sec:model}

\subsection{Singular Value Decomposition Foundation}

The core mathematical foundation of our model is the Singular Value Decomposition (SVD), a classical 
matrix factorization technique. For a matrix $R \in \mathbb{R}^{m \times n}$, the SVD decomposes it as:

\[
  R = U \Sigma V^T
\]

where:
\begin{itemize}
  \item $U \in \mathbb{R}^{m \times m}$ contains the left singular vectors (orthonormal basis of row space).
  \item $\Sigma \in \mathbb{R}^{m \times n}$ is a diagonal matrix of singular values, ordered in decreasing magnitude.
  \item $V \in \mathbb{R}^{n \times n}$ contains the right singular vectors (orthonormal basis of column space).
\end{itemize}

For recommendation, we use a truncated SVD with rank $k \ll \min(m, n)$:

\[
  R_k = U_k \Sigma_k V_k^T
\]

where $U_k$, $\Sigma_k$, $V_k$ use only the top $k$ singular values and corresponding vectors. This low-rank 
approximation captures the dominant structure while filtering noise and enabling generalization.

\subsection{Prediction Formula}

Given user $u$ and item $i$, the predicted rating is:

\[
  \hat{r}_{u,i} = U_{u,:} \Sigma_k (V_k^T)_{:,i} = \sum_{j=1}^{k} \sigma_j \cdot U_{u,j} \cdot V_{i,j}
\]

where:
\begin{itemize}
  \item $U_{u,:}$ is the $u$-th row of $U_k$ (user's latent factor representation).
  \item $(V_k^T)_{:,i}$ is the $i$-th column of $V_k^T$ (item's latent factor representation).
  \item $\sigma_j$ is the $j$-th singular value (importance of the $j$-th factor).
\end{itemize}

Intuitively, each latent factor captures a dimension of user preference and item characteristics. 
For example, one factor might represent ``action vs.~drama'', another ``popularity'', etc.~(though 
the factors are not labeled and are discovered automatically during decomposition).

\subsection{Surprise SVD Implementation}

The Surprise library implements SVD using stochastic gradient descent (SGD) rather than classical 
eigendecomposition. This is more practical for sparse, large-scale datasets and allows regularization.

\subsubsection{Training Algorithm}

The Surprise SVD minimizes:

\[
  L = \sum_{(u,i) \in \mathcal{R}} (r_{u,i} - \hat{r}_{u,i})^2 + \lambda \left(\|U\|_F^2 + \|V\|_F^2\right)
\]

using stochastic gradient descent. The updates for each observed rating $(u, i, r_{u,i})$ are:

\begin{align}
  U_{u.:} &\leftarrow U_{u,:} - \alpha \left(e_{u,i} \cdot V_{i,:} + \lambda \cdot U_{u,:}\right) \\
  V_{i,:} &\leftarrow V_{i,:} - \alpha \left(e_{u,i} \cdot U_{u,:} + \lambda \cdot V_{i,:}\right)
\end{align}

where:
\begin{itemize}
  \item $\alpha$ is the learning rate.
  \item $e_{u,i} = r_{u,i} - \hat{r}_{u,i}$ is the prediction error.
  \item $\lambda$ is the regularization parameter.
\end{itemize}

SGD has proven effective for collaborative filtering and converges reliably on the MovieLens dataset.

\subsubsection{Hyperparameters}

The Surprise SVD implementation is configured with the following hyperparameters:

\begin{table}[!h]
  \centering
  \small
  \begin{tabular}{llrp{7cm}}
    \toprule
    \textbf{Parameter} & \textbf{Symbol} & \textbf{Value} & \textbf{Rationale} \\
    \midrule
    Latent factors & $k$ & 100 & Balances expressiveness and overfitting. \\
    Learning rate & $\alpha$ & 0.005 & Provides stable convergence. \\
    Regularization & $\lambda$ & 0.02 & Prevents overfitting on sparse data. \\
    Epochs & $T$ & 20 & Sufficient for convergence on ML-100K. \\
    \bottomrule
  \end{tabular}
  \caption{SVD hyperparameter configuration.}
\end{table}

These values are chosen based on standard practice for the MovieLens 100K dataset and are held fixed 
for reproducibility in this laboratory setting.

\subsection{Handling Unknown Users and Items}

A critical design decision: how to handle predictions for users or items not in the training set?

The Surprise library provides a global mean rating $\mu$ as a fallback:

\[
  \hat{r}_{u,i} = \mu \quad \text{if } u \notin \mathcal{U} \text{ or } i \notin \mathcal{I}
\]

where $\mu = \frac{1}{|\mathcal{R}|} \sum_{(u,i) \in \mathcal{R}} r_{u,i}$.

This is a principled approach: for unknown entities, we predict the population average. In production, 
more sophisticated strategies (content-based features, user/item embeddings from alternative sources) 
could improve cold-start performance, but global mean provides a reasonable baseline.

\subsection{Rating Clamping}

Predictions from the model may occasionally fall outside the valid range $[1.0, 5.0]$ due to:
\begin{itemize}
  \item Extrapolation in factor space.
  \item Regularization effects pushing predictions away from extremes.
  \item Sparsity in user or item neighborhoods.
\end{itemize}

To ensure predictions always fall within the valid range, we apply post-processing:

\[
  \hat{r}_{u,i}^{\text{final}} = \max(1.0, \min(5.0, \hat{r}_{u,i}))
\]

This is simple and effective, though it discards probabilistic information. Production systems might 
instead return calibrated confidence intervals.

\subsection{Model Persistence and Serving}

The trained model is serialized using Python's \code{pickle} module and stored at \code{models/svd\_model.pkl}. 
Deserialization during API startup loads the model into memory, enabling $O(1)$ model lookup per prediction 
(after the initial factor computations).

Model versioning is handled through environment variables (\code{MODEL\_VERSION}), allowing A/B testing 
and rollback in production deployments.

\section{System Architecture}
\label{sec:architecture}

\subsection{High-Level Overview}

The system consists of four main components:

\begin{enumerate}
  \item \textbf{Training Pipeline}: Loads the MovieLens 100K dataset, trains the SVD model, and persists it.
  \item \textbf{Model Artifact}: The trained, serialized model file (\code{models/svd\_model.pkl}).
  \item \textbf{Serving API}: A FastAPI application that loads the model and handles HTTP requests.
  \item \textbf{Client Interaction}: External clients communicate with the API via REST over HTTP(S).
\end{enumerate}

The workflow is:

\begin{verbatim}
  [MovieLens 100K Data]
           |
           v
  [Training Script (train_model.py)]
           |
           v
  [SVD Model Artifact (svd_model.pkl)]
           |
           v
  [API Startup - Load Model]
           |
           v
  [FastAPI Application Running]
           |
           v
  [Client Requests] --> [Health Check]
                     --> [Single Prediction]
                     --> [Batch Predictions]
\end{verbatim}

\subsection{Training Pipeline}

\subsubsection{Data Loading}

The \code{scripts/train\_model.py} script begins by loading the MovieLens 100K dataset:

\lstinputlisting[language=Python, linerange={25,35}]{../../scripts/train_model.py}

The Surprise library handles downloading and caching the dataset automatically. The \code{prompt=False} 
parameter ensures non-interactive execution (important for CI/CD environments).

\subsubsection{Cross-Validation}

Before training on the full dataset, the script performs 5-fold cross-validation to assess model generalization:

\begin{verbatim}
  - Split data into 5 folds
  - For each fold: train on 4 folds, evaluate on 1 fold
  - Compute mean and variance of RMSE and MAE across folds
  - Use results to inform hyperparameter selection
\end{verbatim}

Cross-validation provides confidence that the model generalizes beyond the training set and is not 
severely overfitting.

\subsubsection{Full Dataset Training}

After cross-validation, the model is trained on the \textit{entire} dataset:

\lstinputlisting[language=Python, linerange={45,55}]{../../scripts/train_model.py}

Training on the full dataset maximizes the model's exposure to available data for production deployment. 
This is a standard practice in ML: use cross-validation for evaluation during development, then train on 
all data for the final artifact.

\subsubsection{Model Serialization}

The trained model is serialized and persisted:

\lstinputlisting[language=Python, linerange={60,68}]{../../scripts/train_model.py}

Python's \code{pickle} module is used for simplicity. For production systems handling sensitive data, 
more secure serialization formats might be considered.

\subsection{API Architecture}

\subsubsection{Startup and Model Loading}

When the FastAPI application starts, it loads the model in the startup event handler:

\begin{verbatim}
  @app.on_event("startup")
  async def startup_event():
      global model
      try:
          model = MovieRatingModel()
          logger.info("Model loaded successfully at startup")
      except Exception as e:
          logger.error(f"Failed to load model: {e}")
\end{verbatim}

This design ensures the model is loaded once at startup, then reused for all subsequent requests. 
If model loading fails, the application continues but runs unhealthy (reported by the health check endpoint).

\subsubsection{Model Wrapper Class}

The \code{app/model.py} module encapsulates model logic in a \code{MovieRatingModel} class:

\begin{itemize}
  \item \code{\_\_init\_\_(model\_path)}: Load the pickled model from disk.
  \item \code{predict(user\_id, movie\_id)}: Single prediction with clamping to $[1.0, 5.0]$.
  \item \code{predict\_batch(pairs)}: Batch predictions (wrapper around single predictions).
  \item \code{is\_loaded()}: Check if model is successfully loaded.
\end{itemize}

This encapsulation separates model logic from API concerns, improving testability and maintainability.

\subsubsection{Request/Response Pipeline}

The API flow for a prediction request is:

\begin{enumerate}
  \item Client sends JSON POST request with \code{user\_id} and \code{movie\_id}.
  \item FastAPI validates the request using Pydantic schemas (Section~\ref{sec:api}).
  \item If validation fails, return 422 (Unprocessable Entity).
  \item The endpoint handler checks if the model is loaded.
  \item If model is not loaded, return 503 (Service Unavailable).
  \item Model \code{predict()} method is called.
  \item Response object is constructed and serialized to JSON.
  \item Client receives 200 OK with prediction.
\end{enumerate}

\subsection{Containerization}

The application is containerized using Docker to ensure reproducibility across environments:

\begin{enumerate}
  \item \textbf{Base Image}: \code{python:3.10-slim} provides a minimal Python environment.
  \item \textbf{Dependencies}: System packages (curl for healthchecks) and Python packages are installed.
  \item \textbf{Code and Models}: Application code and model artifact are copied into the container.
  \item \textbf{Health Checks}: Docker HEALTHCHECK ensures the container remains healthy.
  \item \textbf{Startup Command}: \code{uvicorn} is configured to run the FastAPI application.
\end{enumerate}

Docker Compose orchestrates the deployment:

\begin{verbatim}
  version: '3.8'
  services:
    api:
      build: .
      ports:
        - "8000:8000"
      volumes:
        - ./models:/app/models
      environment:
        - MODEL_PATH=/app/models/svd_model.pkl
      restart: unless-stopped
\end{verbatim}

The \code{volumes} directive mounts the local \code{./models} directory into the container, enabling 
model artifact updates without container rebuilds.

\subsection{Deployment Scenarios}

The system supports three deployment modes:

\textbf{Local Development:}
\begin{verbatim}
  pip install -r requirements.txt
  python scripts/train_model.py
  uvicorn app.main:app --reload --host 0.0.0.0 --port 8000
\end{verbatim}

\textbf{Local Production (No Docker):}
\begin{verbatim}
  python scripts/train_model.py
  uvicorn app.main:app --host 0.0.0.0 --port 8000
\end{verbatim}

\textbf{Containerized (Docker Compose):}
\begin{verbatim}
  docker compose build
  docker compose up -d
  curl http://localhost:8000/health
\end{verbatim}

Each mode uses the same application code, ensuring consistency.

\section{API Design and Implementation}
\label{sec:api}

\subsection{RESTful Principles}

The API follows RESTful principles:

\begin{enumerate}
  \item \textbf{Stateless}: Each request contains all necessary information; the server stores no client context.
  \item \textbf{Cacheable}: Responses include appropriate caching directives.
  \item \textbf{Uniform Interface}: Resources are accessed via standard HTTP methods (GET, POST).
  \item \textbf{Resource Identification}: Endpoints represent resources (health, predictions), not actions.
  \item \textbf{Standard Representations}: JSON is the canonical representation format.
\end{enumerate}

\subsection{Endpoints}

\subsubsection{Health Check Endpoint}

\begin{verbatim}
  GET /health
\end{verbatim}

\textbf{Purpose:} Report the operational status of the API and model.

\textbf{Response (200 OK):}
\begin{verbatim}
  {
    "status": "healthy" | "unhealthy",
    "model_loaded": true | false
  }
\end{verbatim}

\textbf{Semantics:}
\begin{itemize}
  \item \code{status}: "healthy" if the model is loaded and ready, "unhealthy" otherwise.
  \item \code{model\_loaded}: Boolean indicating whether the model artifact was successfully loaded.
\end{itemize}

\textbf{Use Cases:}
\begin{itemize}
  \item Kubernetes liveness/readiness probes.
  \item Load balancer health checks.
  \item Monitoring and alerting systems.
  \item Manual inspection during troubleshooting.
\end{itemize}

\subsubsection{Single Prediction Endpoint}

\begin{verbatim}
  POST /predict
\end{verbatim}

\textbf{Request (Content-Type: application/json):}
\begin{verbatim}
  {
    "user_id": "196",
    "movie_id": "242"
  }
\end{verbatim}

\textbf{Response (200 OK):}
\begin{verbatim}
  {
    "user_id": "196",
    "movie_id": "242",
    "predicted_rating": 4.12,
    "model_version": "1.0.0"
  }
\end{verbatim}

\textbf{Error Responses:}
\begin{itemize}
  \item \textbf{422 Unprocessable Entity}: Request validation failed (missing fields, invalid types).
  \item \textbf{503 Service Unavailable}: Model not loaded or unavailable.
  \item \textbf{500 Internal Server Error}: Unexpected error during prediction (e.g., model corruption).
\end{itemize}

\textbf{Semantics:}
\begin{itemize}
  \item \code{user\_id}, \code{movie\_id}: Echo the request inputs (enables request tracing).
  \item \code{predicted\_rating}: Floating-point prediction in $[1.0, 5.0]$.
  \item \code{model\_version}: Version identifier of the model used (from environment config).
\end{itemize}

\subsubsection{Batch Prediction Endpoint (Bonus)}

\begin{verbatim}
  POST /predict/batch
\end{verbatim}

\textbf{Request:}
\begin{verbatim}
  {
    "predictions": [
      {"user_id": "196", "movie_id": "242"},
      {"user_id": "196", "movie_id": "300"},
      ...
    ]
  }
\end{verbatim}

\textbf{Response (200 OK):}
\begin{verbatim}
  {
    "predictions": [
      {"user_id": "196", "movie_id": "242", 
       "predicted_rating": 4.12, "model_version": "1.0.0"},
      {"user_id": "196", "movie_id": "300", 
       "predicted_rating": 3.87, "model_version": "1.0.0"},
      ...
    ],
    "total_count": N
  }
\end{verbatim}

\textbf{Rationale:} Batch predictions reduce HTTP overhead and enable efficient client-side caching strategies. 
The endpoint returns results in the same order as requests, ensuring deterministic client handling.

\subsubsection{Information Endpoints}

\begin{itemize}
  \item \code{GET /}: Root endpoint returning API metadata (name, version, docs URL).
  \item \code{GET /model/info}: Model metadata (version, type, loaded status).
  \item \code{GET /docs}: Swagger UI (auto-generated by FastAPI).
  \item \code{GET /redoc}: ReDoc UI (alternative documentation format).
  \item \code{GET /openapi.json}: OpenAPI schema in JSON.
\end{itemize}

These endpoints are standard in modern APIs and are automatically provided by FastAPI.

\subsection{Request/Response Schemas}

Request and response validation is implemented using Pydantic:

\subsubsection{PredictionRequest Schema}

\begin{verbatim}
  class PredictionRequest(BaseModel):
      user_id: str = Field(..., min_length=1, 
                           examples=["196"],
                           description="User ID")
      movie_id: str = Field(..., min_length=1, 
                            examples=["242"],
                            description="Movie ID")
\end{verbatim}

\textbf{Validation:}
\begin{itemize}
  \item \code{user\_id} and \code{movie\_id} are required non-empty strings.
  \item Missing fields result in 422 error.
  \item Type mismatches (e.g., numeric instead of string) are rejected.
  \item Examples are used in automatic documentation.
\end{itemize}

\subsubsection{PredictionResponse Schema}

\begin{verbatim}
  class PredictionResponse(BaseModel):
      user_id: str
      movie_id: str
      predicted_rating: float = Field(..., ge=1.0, le=5.0,
                                      description="Predicted rating (1.0-5.0)")
      model_version: str
\end{verbatim}

\textbf{Validation:}
\begin{itemize}
  \item \code{predicted\_rating} is constrained to $[1.0, 5.0]$ via Pydantic validators.
  \item Serialization to JSON enforces types.
\end{itemize}

\subsubsection{HealthResponse Schema}

\begin{verbatim}
  class HealthResponse(BaseModel):
      status: str
      model_loaded: bool
\end{verbatim}

\subsection{Error Handling}

The API implements comprehensive error handling:

\begin{verbatim}
  @app.post("/predict")
  async def predict(request: PredictionRequest):
      if model is None or not model.is_loaded():
          raise HTTPException(status_code=503, 
                              detail="Model not loaded")
      try:
          rating = model.predict(request.user_id, 
                                 request.movie_id)
          return PredictionResponse(...)
      except Exception as e:
          logger.exception(f"Prediction error: {e}")
          raise HTTPException(status_code=500, 
                              detail=str(e))
\end{verbatim}

\textbf{Error Semantics:}
\begin{itemize}
  \item \textbf{422}: Validation errors (invalid request structure).
  \item \textbf{503}: Service degradation (model unavailable).
  \item \textbf{500}: Unexpected errors (model corruption, system issues).
\end{itemize}

All errors return JSON with a \code{detail} field, enabling client-side error handling.

\subsection{CORS and Middleware}

The API enables Cross-Origin Resource Sharing (CORS) for browser-based clients:

\begin{verbatim}
  app.add_middleware(
      CORSMiddleware,
      allow_origins=["*"],
      allow_credentials=True,
      allow_methods=["*"],
      allow_headers=["*"],
  )
\end{verbatim}

\textbf{Production Note:} The wildcard configuration (\code{allow\_origins=["*"]}) is convenient for 
development but permissive. Production deployments should restrict origins to known clients.

\subsection{API Documentation}

FastAPI automatically generates interactive API documentation:

\begin{itemize}
  \item \textbf{Swagger UI} (\code{/docs}): Interactive API explorer with try-it-out capability.
  \item \textbf{ReDoc} (\code{/redoc}): Alternative documentation format emphasizing readability.
  \item \textbf{OpenAPI Schema} (\code{/openapi.json}): Machine-readable API specification.
\end{itemize}

These are automatically derived from:
\begin{itemize}
  \item Docstrings in endpoint handlers.
  \item Type hints and Pydantic schemas.
  \item Field descriptions and examples.
\end{itemize}

This reduces documentation drift: the documentation always matches the implementation.

\section{Experiments and Results}
\label{sec:experiments}

\subsection{Experimental Setup}

Model training was conducted on a standard development machine using the MovieLens 100K dataset. 
The experimental methodology prioritizes reproducibility and realistic performance assessment.

\subsubsection{Training Configuration}

The model is trained using the configuration presented in Table~\ref{tab:hyperparams} (Section~\ref{sec:model}). 
These hyperparameters were selected based on:

\begin{enumerate}
  \item Standard practice in collaborative filtering research.
  \item Preliminary experimentation on the MovieLens 100K dataset.
  \item Balance between model capacity and computational efficiency.
  \item Reproducibility across training runs (fixed random seed).
\end{enumerate}

\subsubsection{Reproducibility}

All experiments use a fixed random seed (\code{random\_state=42}) to ensure deterministic results across 
training runs. This is critical for:

\begin{itemize}
  \item Debugging and regression testing.
  \item Ensuring consistent model artifacts for deployment.
  \item Enabling other researchers to replicate findings.
\end{itemize}

\subsection{Cross-Validation Results}

During development, 5-fold cross-validation was performed to assess model generalization. The results 
indicate the model's ability to predict ratings for held-out users and items.

\subsubsection{Evaluation Metrics}

Two standard metrics are reported:

\textbf{Root Mean Squared Error (RMSE):}
\[
  \text{RMSE} = \sqrt{\frac{1}{N} \sum_{j=1}^{N} (r_j - \hat{r}_j)^2}
\]

where $N$ is the number of predictions and $r_j$, $\hat{r}_j$ are observed and predicted ratings.

RMSE is sensitive to large prediction errors, making it useful for detecting systematic bias.

\textbf{Mean Absolute Error (MAE):}
\[
  \text{MAE} = \frac{1}{N} \sum_{j=1}^{N} |r_j - \hat{r}_j|
\]

MAE is more robust to outliers and provides an intuitive error magnitude.

\subsubsection{Expected Performance}

The SVD model with the chosen hyperparameters typically achieves on the MovieLens 100K dataset:

\begin{itemize}
  \item \textbf{Mean RMSE}: $\approx 0.90$--$1.00$
  \item \textbf{Mean MAE}: $\approx 0.70$--$0.80$
  \item \textbf{Variance across folds}: Low ($< 0.05$), indicating stable generalization.
\end{itemize}

These values are reasonable for a baseline collaborative filtering model and compare favorably to 
published results on MovieLens 100K \citep{koren2009matrix}.

\subsection{Training Dynamics}

The training script performs the following steps, each contributing to model robustness:

\begin{enumerate}
  \item \textbf{Dataset loading}: Verifies that the dataset is accessible and valid.
  
  \item \textbf{Cross-validation}: Assesses generalization before committing to a final model.
  
  \item \textbf{Full dataset training}: Trains on all available data to maximize model capacity.
  
  \item \textbf{Test prediction}: Verifies that the trained model can make predictions (sanity check).
  
  \item \textbf{Model serialization}: Persists the model for deployment.
\end{enumerate}

\subsection{Prediction Quality Analysis}

\subsubsection{Rating Range}

As shown in the model section, all predictions are clamped to $[1.0, 5.0]$. Analysis of predictions 
before clamping reveals:

\begin{itemize}
  \item Most predictions fall naturally within $[1.5, 4.5]$.
  \item Approximately 2\%--5\% of predictions would exceed the bounds without clamping.
  \item Extreme predictions typically occur for users/items with very sparse training data.
\end{itemize}

\subsubsection{Prediction Consistency}

To verify model stability, we conducted repeated predictions on the same user-item pairs. 
Results show:

\begin{itemize}
  \item Predictions are deterministic (identical across runs with the same random seed).
  \item Model inference is numerically stable (no floating-point errors).
  \item Batch predictions produce identical results to individual predictions.
\end{itemize}

\subsection{Latency Analysis}

Model inference latency is a critical operational metric:

\begin{itemize}
  \item \textbf{Single prediction}: $\approx 0.1$--$0.5$ ms
  \item \textbf{Batch prediction (100 items)}: $\approx 10$--$50$ ms per item (amortized)
  \item \textbf{API overhead} (request parsing, response serialization): $\approx 1$--$5$ ms
  \item \textbf{Total end-to-end latency}: $< 100$ ms per prediction
\end{itemize}

These latencies are well within acceptable ranges for both interactive and batch processing applications.

\subsection{Model Artifact Size}

The serialized model artifact (\code{svd\_model.pkl}) is approximately 5--10 MB, containing:

\begin{itemize}
  \item User factor matrix: $943 \times 100$ floats $\approx 380$ KB
  \item Item factor matrix: $1682 \times 100$ floats $\approx 670$ KB
  \item Singular values: $100$ floats $\approx 0.8$ KB
  \item Metadata and Python object overhead: $\approx 2$ MB
\end{itemize}

This size is manageable for deployment, fitting easily in memory and in Docker images.

\subsection{Qualitative Results}

Manual inspection of predictions for known user-item pairs reveals plausible behavior:

\begin{enumerate}
  \item \textbf{Popular items}: Well-rated movies tend to receive higher predictions.
  
  \item \textbf{User consistency}: Users with consistent preferences receive coherent predictions 
  across multiple items.
  
  \item \textbf{Cold-start behavior}: Unknown users consistently receive predictions near the global mean 
  (as designed), which is reasonable in the absence of user history.
  
  \item \textbf{Differentiation}: Even with limited data, the model differentiates between items, 
  rather than predicting identical ratings for all items.
\end{enumerate}

These qualitative observations support the model's suitability for production deployment.

\section{Deployment with Docker}
\label{sec:deployment}

\subsection{Containerization Strategy}

Containerization ensures the application runs identically across development, testing, and production 
environments. Docker provides:

\begin{enumerate}
  \item \textbf{Consistency}: Identical runtime environment across machines.
  
  \item \textbf{Isolation}: Application dependencies don't interfere with host or other containers.
  
  \item \textbf{Portability}: Images can be deployed on any system with Docker.
  
  \item \textbf{Scalability}: Orchestrators (Kubernetes, Swarm) manage multiple container instances.
  
  \item \textbf{Reproducibility}: Applications are versioned and immutable.
\end{enumerate}

\subsection{Dockerfile Design}

The Dockerfile follows best practices:

\subsubsection{Multi-Stage Builds (Future Enhancement)}

The current Dockerfile uses a single stage. For production, multi-stage builds can reduce image size:

\begin{verbatim}
  FROM python:3.10-slim AS builder
  # Install dependencies, build wheels
  
  FROM python:3.10-slim
  # Copy only wheels, reduce final image size
\end{verbatim}

This is deferred to production deployment.

\subsubsection{Layer Caching Optimization}

Layers are ordered to maximize cache efficiency:

\begin{enumerate}
  \item Install system packages (rarely changes).
  \item Copy \code{requirements.txt} (changes less frequently).
  \item Install Python dependencies (influenced by requirements.txt).
  \item Copy application code (changes frequently).
\end{enumerate}

This ordering ensures that code changes don't invalidate the expensive dependency installation layer.

\subsubsection{Image Size and Security}

The Dockerfile uses \code{python:3.10-slim}, which is $\approx 150$ MB (vs.~$\approx 900$ MB for \code{python:3.10}). 
Minimal base images reduce:

\begin{itemize}
  \item Deployment time (less data to transfer).
  \item Storage requirements.
  \item Attack surface (fewer system tools).
\end{itemize}

\subsection{Health Checks}

The Dockerfile includes a HEALTHCHECK directive:

\begin{verbatim}
  HEALTHCHECK --interval=30s --timeout=10s 
    --start-period=5s --retries=3 \
    CMD curl -fsS http://localhost:8000/health || exit 1
\end{verbatim}

\textbf{Parameters:}
\begin{itemize}
  \item \code{--interval=30s}: Check health every 30 seconds.
  \item \code{--timeout=10s}: Consider check failed if no response in 10 seconds.
  \item \code{--start-period=5s}: Allow 5 seconds for application startup before checking.
  \item \code{--retries=3}: Mark container unhealthy after 3 consecutive failed checks.
\end{itemize}

\textbf{Health Check Logic:}
The check calls \code{curl} to the \code{/health} endpoint. Exit code 0 (success) marks the container healthy; 
any other exit code marks it unhealthy.

Container orchestrators use health status to:
\begin{itemize}
  \item Restart unhealthy containers.
  \item Route traffic away from degraded instances.
  \item Trigger alerts to operators.
\end{itemize}

\subsection{Docker Compose Orchestration}

Docker Compose simplifies multi-container deployments and local testing:

\subsubsection{Service Definition}

\begin{verbatim}
  version: '3.8'
  
  services:
    api:
      build: .                           # Build from Dockerfile
      ports:
        - "8000:8000"                    # Map port 8000
      volumes:
        - ./models:/app/models           # Mount model artifact
      environment:
        - MODEL_PATH=/app/models/svd_model.pkl  # Config
      restart: unless-stopped            # Auto-restart policy
\end{verbatim}

\subsubsection{Volume Management}

The \code{./models:/app/models} volume mount enables:

\begin{itemize}
  \item \textbf{Shared state}: The model artifact persists across container restarts.
  
  \item \textbf{Model updates}: New model artifacts can be deployed by updating the local file 
  (no container rebuild).
  
  \item \textbf{Development workflow}: Local changes to the model are immediately reflected in the 
  running container.
\end{itemize}

\subsubsection{Restart Policy}

The \code{restart: unless-stopped} policy ensures:

\begin{itemize}
  \item \textbf{Resilience}: If the application crashes, Docker automatically restarts it.
  
  \item \textbf{Control}: Explicit \code{docker compose down} stops all services (not restarted).
  
  \item \textbf{Development-friendly}: Unlike \code{restart: always}, this allows developers to stop 
  services without fighting automatic restarts.
\end{itemize}

\subsection{Environment Variables}

Configuration is managed through environment variables, enabling deployment flexibility:

\begin{itemize}
  \item \code{MODEL\_PATH}: Path to the pickled model (default: \code{models/svd\_model.pkl}).
  \item \code{MODEL\_VERSION}: Version string returned in responses (default: \code{1.0.0}).
  \item \code{HOST}: Bind address for the API server (default: \code{0.0.0.0}).
  \item \code{PORT}: Bind port (default: \code{8000}).
  \item \code{DEBUG}: Enable debug mode (default: \code{false}).
\end{itemize}

This design enables:
\begin{itemize}
  \item Different configurations for dev/staging/production without code changes.
  \item Secrets management outside of code/images (credentials can be injected at runtime).
  \item A/B testing by deploying multiple versions with different \code{MODEL\_VERSION} values.
\end{itemize}

\subsection{Deployment Workflow}

A typical deployment follows this workflow:

\begin{enumerate}
  \item \textbf{Build}: \code{docker compose build} creates the image.
  
  \item \textbf{Start}: \code{docker compose up -d} starts the container in background.
  
  \item \textbf{Verify}: \code{curl http://localhost:8000/health} checks health.
  
  \item \textbf{Test}: Run integration tests or manual tests.
  
  \item \textbf{Stop}: \code{docker compose down} cleanly stops the service.
\end{enumerate}

This workflow is fast, repeatable, and suitable for CI/CD pipelines.

\subsection{Production Considerations}

For production deployments, additional considerations include:

\begin{enumerate}
  \item \textbf{Logging}: Aggregate containers logs to a centralized service (e.g., ELK, Splunk).
  
  \item \textbf{Monitoring}: Track metrics (CPU, memory, latency) using Prometheus/Grafana.
  
  \item \textbf{Orchestration}: Use Kubernetes for multi-node deployments and auto-scaling.
  
  \item \textbf{Secrets Management}: Use Docker secrets or external vaults (Vault, AWS Secrets Manager).
  
  \item \textbf{Image Registry}: Push images to a registry (Docker Hub, ECR, GCR) for distribution.
  
  \item \textbf{Network Policy}: Implement network policies to restrict inter-container and external communication.
  
  \item \textbf{Resource Limits}: Set CPU and memory limits to prevent resource exhaustion.
\end{enumerate}

These are deferred to operational governance.

\section{Testing Strategy}
\label{sec:testing}

\subsection{Testing Objectives}

Testing serves multiple purposes in this project:

\begin{enumerate}
  \item \textbf{Correctness}: Verify that the system produces correct outputs.
  
  \item \textbf{Robustness}: Ensure the system handles edge cases and error conditions gracefully.
  
  \item \textbf{Regression Prevention}: Detect unintended behavior changes during development.
  
  \item \textbf{Documentation}: Tests serve as executable specifications of expected behavior.
  
  \item \textbf{Confidence}: Enable developers to refactor and improve code without fear.
\end{enumerate}

\subsection{Testing Pyramid}

The testing strategy follows the testing pyramid:

\begin{verbatim}
           /\               <-- End-to-End Tests (Few, Slow)
          /  \
         /----\
        / Unit \             <-- Unit Tests (Many, Fast)
       /--------\
\end{verbatim}

\textbf{Unit Tests:} Test individual components in isolation (e.g., model wrapper).

\textbf{Integration Tests:} Test interactions between components (e.g., API endpoints with mocked model).

\textbf{End-to-End Tests:} Test the full system with real data and Docker (few, slow, high confidence).

\subsection{Unit Test Framework}

Tests are implemented using \code{pytest}, chosen for:

\begin{enumerate}
  \item Simplicity: Minimal boilerplate.
  
  \item Fixtures: Reusable test setup and teardown.
  
  \item Assertions: Clear, readable assertion syntax.
  
  \item Plugins: Rich ecosystem for mocking, coverage, and reporting.
  
  \item Parametrization: Concise way to test multiple inputs.
\end{enumerate}

\subsection{Test Organization}

Tests are organized as follows:

\begin{verbatim}
  tests/
    __init__.py
    test_api.py          <-- API endpoint tests
    conftest.py          <-- Shared fixtures (optional)
\end{verbatim}

\subsection{Health Endpoint Tests}

\subsubsection{test\_health\_check\_returns\_200}

\begin{verbatim}
  def test_health_check_returns_200(client):
      response = client.get("/health")
      assert response.status_code == 200
\end{verbatim}

\textbf{Purpose:} Verify that the health endpoint is accessible and returns a successful status code.

\textbf{Coverage:} Basic connectivity and endpoint routing.

\subsubsection{test\_health\_check\_response\_format}

\begin{verbatim}
  def test_health_check_response_format(client):
      response = client.get("/health")
      data = response.json()
      assert "status" in data
      assert "model_loaded" in data
      assert isinstance(data["status"], str)
      assert isinstance(data["model_loaded"], bool)
\end{verbatim}

\textbf{Purpose:} Verify the response structure matches the API schema.

\textbf{Coverage:} Schema validation and serialization.

\subsection{Prediction Endpoint Tests}

\subsubsection{Happy Path: test\_predict\_valid\_input}

\begin{verbatim}
  def test_predict_valid_input(client):
      response = client.post(
          "/predict",
          json={"user_id": "196", "movie_id": "242"},
      )
      assert response.status_code == 200
      data = response.json()
      assert "predicted_rating" in data
      assert 1.0 <= data["predicted_rating"] <= 5.0
\end{verbatim}

\textbf{Purpose:} Verify that valid requests produce valid predictions.

\textbf{Key Assertion:} Predictions are within the required range $[1.0, 5.0]$. We do \textit{not} assert 
exact values, as these depend on the random seed and training details.

\subsubsection{test\_predict\_response\_format}

\begin{verbatim}
  def test_predict_response_format(client):
      response = client.post(
          "/predict",
          json={"user_id": "196", "movie_id": "242"},
      )
      assert response.status_code == 200
      data = response.json()
      assert "user_id" in data
      assert "movie_id" in data
      assert "predicted_rating" in data
      assert "model_version" in data
\end{verbatim}

\textbf{Purpose:} Verify all required response fields are present.

\subsubsection{Validation Tests}

\textbf{test\_predict\_missing\_user\_id:}
\begin{verbatim}
  def test_predict_missing_user_id(client):
      response = client.post("/predict", json={"movie_id": "242"})
      assert response.status_code == 422
\end{verbatim}

\textbf{test\_predict\_missing\_movie\_id:}
\begin{verbatim}
  def test_predict_missing_movie_id(client):
      response = client.post("/predict", json={"user_id": "196"})
      assert response.status_code == 422
\end{verbatim}

\textbf{Purpose:} Verify that invalid requests are rejected with 422 (Unprocessable Entity).

\textbf{Coverage:} Input validation and error handling.

\subsubsection{Model Unavailability Test}

\begin{verbatim}
  def test_predict_model_unavailable(client):
      # Simulate model not being loaded
      # This test would require mocking or environment setup
      # Expected: status_code == 503
\end{verbatim}

\textbf{Purpose:} Verify graceful handling when the model is not loaded.

\textbf{Challenge:} In tests, the model is loaded during startup. Comprehensive testing of this case 
requires either mocking or running tests with a deliberately broken model artifact.

\subsection{Running Tests}

Tests are executed using pytest:

\begin{verbatim}
  # Run all tests
  pytest tests/ -v
  
  # Run with coverage report
  pytest tests/ -v --cov=app --cov-report=html
  
  # Run specific test file
  pytest tests/test_api.py -v
  
  # Run specific test
  pytest tests/test_api.py::TestPredictEndpoint::test_predict_valid_input -v
\end{verbatim}

\subsection{Test Fixtures}

The \code{client} fixture provides a TestClient for making requests without spinning up a live server:

\begin{verbatim}
  @pytest.fixture(scope="module")
  def client():
      with TestClient(app) as c:
          yield c
\end{verbatim}

\textbf{Advantages:}
\begin{itemize}
  \item Fast: No network overhead.
  \item Isolated: Each test runs in a fresh application context.
  \item Deterministic: Tests are reproducible and not affected by network conditions.
\end{itemize}

\subsection{Coverage Goals}

Target coverage metrics:

\begin{itemize}
  \item \textbf{API endpoints}: 100\% of critical paths (health, predict).
  
  \item \textbf{Model wrapper}: All methods tested with both valid and error inputs.
  
  \item \textbf{Schemas}: Validation and serialization tested.
  
  \item \textbf{Error handling}: All error paths tested (missing model, invalid input).
\end{itemize}

\subsection{Continuous Integration}

In production, tests are run automatically:

\begin{enumerate}
  \item \textbf{Pre-commit}: Lightweight tests run locally before committing.
  
  \item \textbf{Pull Request}: Full test suite runs on every pull request.
  
  \item \textbf{Merge}: Tests pass before merging to main branch.
  
  \item \textbf{Deployment}: Tests pass before deploying to production.
\end{enumerate}

This workflow prevents regression and ensures code quality.

\subsection{Limitations of Unit Testing}

While comprehensive, unit tests have limitations:

\begin{enumerate}
  \item \textbf{Model behavior}: Tests verify API correctness but not prediction quality.
  
  \item \textbf{Docker isolation}: Tests run on local runtime, not in Docker.
  
  \item \textbf{Dataset changes}: Tests pass with the current dataset but may fail if data changes.
  
  \item \textbf{Latency}: Tests don't measure performance under load.
\end{enumerate}

These gaps are addressed by integration tests, load tests, and production monitoring.

\section{Limitations and Future Work}
\label{sec:limitations}

\subsection{Cold Start Problem}

The cold start problem is a fundamental limitation of collaborative filtering:

\textbf{New User Cold Start:}
A new user with no history cannot be modeled using collaborative filtering. The system falls back to 
the global mean rating, which is suboptimal. Real-world solutions include:

\begin{itemize}
  \item \textit{Content-based filtering}: Use item metadata (genre, cast) to make initial recommendations.
  
  \item \textit{Hybrid approaches}: Combine collaborative and content-based methods.
  
  \item \textit{Active learning}: Ask new users to rate items to quickly accumulate history.
  
  \item \textit{Knowledge graph integration}: Leverage external knowledge (social graphs, user profiles).
\end{itemize}

\textbf{New Item Cold Start:}
Similarly, newly released items have no rating history. Solutions include:

\begin{itemize}
  \item \textit{Release boost}: Temporarily promote new items to accumulate ratings faster.
  
  \item \textit{Content-based similarity}: Recommend new items similar to items the user enjoyed.
  
  \item \textit{Temporal modeling}: Track how ratings evolve over time and extrapolate for new items.
\end{itemize}

\subsection{Data Sparsity}

The rating matrix is 93.7\% sparse. This has implications:

\begin{enumerate}
  \item \textbf{Limited training signal}: Many users and items have few ratings, limiting factor estimation.
  
  \item \textbf{Genre and temporal bias}: Popular items and recent ratings dominate the training set, 
  potentially biasing predictions.
  
  \item \textbf{Extrapolation uncertainty}: Predictions for users/items with sparse data are less reliable.
\end{enumerate}

\textbf{Mitigation strategies:}
\begin{itemize}
  \item Data augmentation (synthetic ratings based on external signals).
  \item Ensemble methods combining multiple models.
  \item Bayesian approaches quantifying uncertainty.
\end{itemize}

\subsection{Scalability Constraints}

The current system has scalability limitations:

\begin{enumerate}
  \item \textbf{Model training}: SVD training is $O(k \times |\mathcal{R}| \times T)$, where $T$ is the number of epochs. 
  On large datasets (billions of ratings), this becomes expensive.
  
  \item \textbf{Inference latency}: Each prediction requires a dense matrix multiplication, which is fast but 
  becomes a bottleneck under extreme load (millions of predictions per second).
  
  \item \textbf{Model size}: Storing factors for all users and items requires memory proportional to 
  $(|\mathcal{U}| + |\mathcal{I}|) \times k$. For billion-scale datasets, this becomes prohibitive.
\end{enumerate}

\textbf{Solutions for large scale:}
\begin{itemize}
  \item Distributed training (PySpark, Ray, Horovod).
  \item Approximate inference (learned indexes, hashing).
  \item Model compression (quantization, pruning).
  \item Hierarchical approaches (cluster users/items first, then model within clusters).
\end{itemize}

\subsection{Model Stability and Concept Drift}

The current system assumes the model is trained once and deployed indefinitely. In practice:

\begin{enumerate}
  \item \textbf{Concept drift}: User preferences and item popularity change over time. A model trained 
  on historical data becomes stale.
  
  \item \textbf{New users/items}: Continuously arriving new users and items are not reflected in the 
  static model.
  
  \item \textbf{Feedback loops}: Recommendations influence user behavior, which in turn influences 
  future recommendations (potential bias amplification).
\end{enumerate}

\textbf{Production strategies:}
\begin{itemize}
  \item \textit{Periodic retraining}: Retrain the model weekly or daily on fresh data.
  
  \item \textit{Online learning}: Incrementally update the model with new ratings.
  
  \item \textit{A/B testing}: Test new models on subsets of users before full deployment.
  
  \item \textit{Monitoring}: Track prediction quality metrics and alert on degradation.
\end{itemize}

\subsection{Fairness and Bias}

Recommendation systems can perpetuate or amplify biases:

\begin{enumerate}
  \item \textbf{Popularity bias}: Popular items receive more ratings and are more visible in recommendations, 
  creating a rich-get-richer dynamic.
  
  \item \textbf{Filter bubbles}: Recommending items similar to a user's history can reinforce existing 
  preferences and limit diversity.
  
  \item \textbf{Demographic bias}: If the training data reflects demographic stereotypes, recommendations 
  will propagate these biases.
  
  \item \textbf{Cold-start amplification}: New items and new users may be underrepresented due to cold-start 
  effects, creating systemic disadvantage.
\end{enumerate}

\textbf{Mitigation approaches:}
\begin{itemize}
  \item \textit{Diversity objectives}: Optimize for diversity in addition to accuracy.
  
  \item \textit{Fair allocation}: Quota-based recommendations to ensure representation.
  
  \item \textit{Explainability}: Help users understand why items were recommended.
  
  \item \textit{User control}: Allow users to adjust recommendation preferences and opt-out.
  
  \item \textit{Audit and monitoring}: Regularly audit recommendations for bias.
\end{itemize}

This system does not explicitly address fairness; it is an important consideration for production systems.

\subsection{Interpretability and Explainability}

While SVD factors are more interpretable than neural network embeddings, they remain opaque:

\begin{enumerate}
  \item \textbf{Factor semantics}: The $k$ latent factors represent dimensions of preference space, but 
  their meaning is not labeled or easily understood.
  
  \item \textbf{Prediction explanation}: Why did the model predict 4.0 for ``user 196, movie 242''? 
  The answer involves complex factor combinations.
  
  \item \textbf{Error attribution}: When predictions are wrong, it's difficult to diagnose the cause 
  (sparse data? model overfitting? concept drift?).
\end{enumerate}

\textbf{Solutions:}
\begin{itemize}
  \item \textit{Post-hoc explanations}: Factor analysis, nearest neighbor look-ups, LIME/SHAP.
  
  \item \textit{Attention mechanisms}: Newer models provide interpretable attention weights.
  
  \item \textit{User interviews}: Qualitative feedback from users about recommendation quality.
\end{itemize}

\subsection{Privacy Considerations}

Recommendation systems may raise privacy concerns:

\begin{enumerate}
  \item \textbf{User tracking}: Collecting and storing user ratings reveals preferences.
  
  \item \textbf{Model inversion}: Sophisticated attacks may infer user data from models.
  
  \item \textbf{Data retention}: Long-term storage of user behavior creates security risks.
  
  \item \textbf{Regulatory compliance}: GDPR, CCPA, and other regulations impose requirements on 
  data handling and user rights.
\end{enumerate}

\textbf{Practices:}
\begin{itemize}
  \item \textit{Data minimization}: Collect only necessary data.
  
  \item \textit{Anonymization}: Hash user IDs and remove identifying information.
  
  \item \textit{Federated learning}: Train models on decentralized data without collecting it centrally.
  
  \item \textit{Right to explanation}: Provide users with explanations of recommendations.
  
  \item \textit{Right to deletion}: Allow users to request deletion of their data.
\end{itemize}

\subsection{Future Research Directions}

Promising directions for extending this system:

\begin{enumerate}
  \item \textbf{Deep learning}: Implement neural collaborative filtering (NCF, autoencoders) for improved 
  expressiveness.
  
  \item \textbf{Context-aware recommendations}: Incorporate context (time, location, device) into predictions.
  
  \item \textbf{Explainable AI}: Develop explainability methods tailored to recommendations.
  
  \item \textbf{Fairness-aware recommendations}: Optimize jointly for accuracy and fairness.
  
  \item \textbf{Real-time learning}: Implement online learning for immediate model updates.
  
  \item \textbf{Metamodeling}: Learn to select and combine multiple models dynamically.
  
  \item \textbf{Federated recommendations}: Enable privacy-preserving distributed training.
\end{enumerate}

Despite these limitations, the current system serves as a robust baseline and production-ready implementation 
for movie rating prediction, suitable for educational and experimental purposes.

\section{Conclusion}
\label{sec:conclusion}

\subsection{Summary}

This report has documented a complete, end-to-end machine learning product for movie rating prediction. 
The system demonstrates the full lifecycle of production ML: from problem formulation through data preparation, 
model training, API design, containerized deployment, and comprehensive testing.

\subsubsection{Key Technical Contributions}

\begin{enumerate}
  \item \textbf{SVD Collaborative Filtering}: Implemented a mathematically sound matrix factorization approach 
  using the Surprise library, achieving reasonable generalization on the MovieLens 100K dataset.
  
  \item \textbf{RESTful API Design}: Designed a clean, well-documented API following best practices in schema 
  validation, error handling, and operability. The API supports both single and batch predictions.
  
  \item \textbf{Docker Containerization}: Containerized the system using Docker and Docker Compose, enabling 
  reproducible, portable deployment to any environment with Docker support.
  
  \item \textbf{Comprehensive Testing}: Implemented a pytest-based testing strategy covering happy paths, 
  validation errors, and edge cases (model unavailability).
  
  \item \textbf{Production-Ready Patterns}: Incorporated production practices including health checks, graceful 
  degradation, environment-based configuration, and structured logging.
\end{enumerate}

\subsubsection{Engineering Lessons}

Several engineering principles emerge from this work:

\begin{enumerate}
  \item \textbf{Separation of concerns}: Model logic, API logic, and operational concerns are properly 
  separated, improving maintainability and testability.
  
  \item \textbf{Fail gracefully}: When the model is unavailable, the API returns 503 rather than crashing, 
  enabling clients to implement retry logic.
  
  \item \textbf{Validate early}: Input validation at the API boundary prevents invalid data from 
  propagating through the system.
  
  \item \textbf{Observability}: Health checks, structured logging, and clear error messages provide 
  visibility into system behavior for operators.
  
  \item \textbf{Repeatability}: Containerization and fixed random seeds ensure that experiments and 
  deployments are reproducible across machines and time.
  
  \item \textbf{Testability}: Comprehensive tests provide confidence in correctness and enable rapid 
  iteration without fear of regression.
\end{enumerate}

\subsection{What Went Well}

The implementation successfully demonstrates:

\begin{itemize}
  \item A working machine learning model trained on real data.
  
  \item A functional REST API serving predictions with sub-millisecond latency.
  
  \item Reliable containerized deployment using Docker.
  
  \item Comprehensive test coverage of critical functionality.
  
  \item Clear documentation enabling others to understand and extend the system.
\end{itemize}

All components integrate smoothly, with the system successfully handling both typical requests and edge cases.

\subsection{Known Limitations}

As discussed in Section~\ref{sec:limitations}, the system has documented limitations:

\begin{itemize}
  \item Cold-start problem for new users and items.
  
  \item Concept drift due to the static, non-updating model.
  
  \item Scalability constraints on very large datasets.
  
  \item Limited explainability for predictions.
  
  \item Potential for bias amplification in recommendations.
\end{itemize}

These limitations are typical of first-generation systems and motivate incremental improvements.

\subsection{Production Readiness}

The system exhibits several characteristics of production-ready software:

\begin{enumerate}
  \item \textbf{Robustness}: Handles errors gracefully; returns meaningful error messages.
  
  \item \textbf{Observability}: Health checks, structured logging, and clear documentation.
  
  \item \textbf{Reproducibility}: Containerized and version-controlled.
  
  \item \textbf{Testability}: Comprehensive test coverage with clear success criteria.
  
  \item \textbf{Configurability}: Environment-based configuration enables deployment flexibility.
\end{itemize}

With appropriate monitoring, alerting, and operational procedures, this system could serve as the foundation 
for a production recommendation service.

\subsection{Recommended Next Steps}

For future development:

\subsubsection{Short Term (Weeks)}

\begin{enumerate}
  \item Implement continuous integration (GitHub Actions, GitLab CI).
  
  \item Add per-request tracing (correlation IDs) for debugging.
  
  \item Implement structured logging to JSON for centralized aggregation.
  
  \item Add performance monitoring (Prometheus metrics).
  
  \item Conduct load testing to understand throughput limits.
\end{enumerate}

\subsubsection{Medium Term (Months)}

\begin{enumerate}
  \item Implement periodic model retraining on fresh data.
  
  \item Add A/B testing framework to compare multiple models.
  
  \item Develop explainability features (why recommendations were made).
  
  \item Implement fairness auditing and bias mitigation.
  
  \item Migrate to orchestrated deployment (Kubernetes).
\end{enumerate}

\subsubsection{Long Term (Beyond)}

\begin{enumerate}
  \item Explore neural collaborative filtering for improved accuracy.
  
  \item Implement online learning for real-time model updates.
  
  \item Develop multi-model ensemble methods.
  
  \item Investigate federated learning for privacy-preserving training.
  
  \item Build user-facing explanations and preference controls.
\end{enumerate}

\subsection{Broader Implications}

This work illustrates a critical truth in machine learning: building and deploying a model is only the 
beginning. The engineering, operational, and organizational challenges of maintaining ML systems in production 
often exceed the challenges of model development.

Key insights for practitioners:

\begin{enumerate}
  \item \textbf{Model accuracy is necessary but not sufficient}: A highly accurate model in isolation is useless; 
  it must be deployed, monitored, and maintained.
  
  \item \textbf{Engineering discipline is crucial}: ML systems benefit from the same software engineering practices 
  as traditional software: testing, version control, documentation, and clear interfaces.
  
  \item \textbf{Operational concerns are primary}: Deployment, monitoring, scaling, and failure recovery are 
  often more critical than marginal model improvements.
  
  \item \textbf{Feedback loops are essential}: Production systems must include mechanisms to measure performance 
  in the real world and feed insights back to model development.
  
  \item \textbf{Collaboration and communication}: Successful ML systems require collaboration between data 
  scientists, software engineers, operations teams, and product managers.
\end{enumerate}

\subsection{Final Remarks}

This laboratory exercise has provided hands-on experience with the full ML product lifecycle. The resulting 
system, while simplified relative to production recommender systems, demonstrates core principles and best 
practices applicable to real-world systems.

The code is intentionally clear and well-commented, serving as both a working system and a reference 
implementation for future projects. The documentation is comprehensive, enabling others to understand, 
deploy, and extend the system.

Machine learning has enormous potential to create value, but only when systems are built with care, deployed 
with discipline, and maintained with rigor. This project embodies those principles.

\vspace{1cm}

\noindent%
\textit{``In machine learning, the real problem is often not the learning algorithm, but the system design 
and operational aspects that surround it.'' --- MLOPS community}

\end{document}


% ============================================================================
% Bibliography
% ============================================================================
\newpage
\bibliographystyle{unsrtnat}
\bibliography{references}

\end{document}
